\documentclass[11pt]{article}
\usepackage[margin=1in]{geometry}
\usepackage{amsmath,amssymb,amsthm}
\usepackage{array}
\usepackage{parskip}
\usepackage{booktabs}
\setlength{\emergencystretch}{3em}

\newtheorem{theorem}{Theorem}
\newtheorem{lemma}[theorem]{Lemma}
\newtheorem{corollary}[theorem]{Corollary}
\newtheorem{proposition}[theorem]{Proposition}
\theoremstyle{definition}
\newtheorem{definition}[theorem]{Definition}
\theoremstyle{remark}
\newtheorem{remark}[theorem]{Remark}

\newcommand{\N}{\mathbb{N}}
\newcommand{\Z}{\mathbb{Z}}
\newcommand{\Wtwelve}{W_{12}}

\title{Disproof of conjecture \texttt{00000008417}}
\author{earthking11}
\date{2026-09-14}

\begin{document}
\maketitle

\section*{Verdict: FALSE}

We disprove conjecture \texttt{00000008417} as stated. The witness is the small
Witt design $\Wtwelve = S(5,6,12)$: a genuine $5$-$(12,6,1)$ design with $132$
blocks. It exists \emph{at} $n = 12$, and at $(t,k,n) = (5,6,12)$ all six
classical divisibility conditions hold. Hence the existence threshold obeys
$n_0(5,6) \le 12$, which contradicts the filed conjunct ``the lower bound
$n_0(5,6) > 12$ is given by a concrete divisibility obstruction''. In fact there
is no divisibility obstruction at $n = 12$ at all. The refutation is
unconditional and completely elementary; it does not rely on any unproved
hypothesis.

\section{The conjecture, quoted verbatim}

From \texttt{conjectures/00000008417.md}:

\begin{quote}
\textbf{English.} Definition: A Steiner system $S(t,k,n)$ is a design in which
every $t$-element subset lies in exactly one block. Conjecture: $S(t,k,n)$
exists if and only if the classical divisibility conditions hold and $n$ is at
least $n_0(t,k)$; $n_0(t,k)$ admits an explicit bound of type $(kt)^{ct}$, and
the lower bound $n_0(5,6) > 12$ is given by a concrete divisibility obstruction.
(explicit threshold for Steiner systems)
\end{quote}

The conjecture is filed in both English and Chinese; the Chinese original is
reproduced verbatim (in Unicode) in \texttt{README.md}. The PDF font does not
embed CJK glyphs, so the Chinese is kept in the markdown copy rather than
duplicated here.

We read the conjecture as the conjunction of two clauses:
\begin{enumerate}
\item[(C1)] \emph{(the existence criterion)} for every $(t,k,n)$ with
  $t < k \le n$, $S(t,k,n)$ exists if and only if the classical divisibility
  conditions hold and $n \ge n_0(t,k)$;
\item[(C2)] \emph{(the small-$n$ lower bound)} $n_0(5,6) > 12$, and this bound
  is ``given by a concrete divisibility obstruction''.
\end{enumerate}
We show that (C1) and (C2) cannot both be true, and that (C2) is false as filed.

\section{The divisibility conditions at $(5,6,12)$}

For an $S(t,k,n)$ the classical divisibility conditions are
\[
  \lambda_i \;:=\; \frac{\binom{n-i}{t-i}}{\binom{k-i}{t-i}} \;\in\; \Z ,
  \qquad i = 0,1,\dots,t ,
\]
where $\lambda_i$ is the number of blocks through a fixed $i$-subset ($\lambda_0$
is the total number of blocks). At $(t,k,n) = (5,6,12)$:
\[
  \lambda_0 = \frac{\binom{12}{5}}{\binom{6}{5}} = \frac{792}{6} = 132,
  \quad
  \lambda_1 = \frac{\binom{11}{4}}{\binom{5}{4}} = \frac{330}{5} = 66,
  \quad
  \lambda_2 = \frac{\binom{10}{3}}{\binom{4}{3}} = \frac{120}{4} = 30,
\]
\[
  \lambda_3 = \frac{\binom{9}{2}}{\binom{3}{2}} = \frac{36}{3} = 12,
  \quad
  \lambda_4 = \frac{\binom{8}{1}}{\binom{2}{1}} = \frac{8}{2} = 4,
  \quad
  \lambda_5 = \frac{\binom{7}{0}}{\binom{1}{0}} = \frac{1}{1} = 1 .
\]
Every quotient is an integer. We record this as exact products:
\begin{center}
\begin{tabular}{r r r}
\toprule
$i$ & exact identity & quotient $\lambda_i$ \\
\midrule
$0$ & $792 = 132 \cdot 6$ & $132$ \\
$1$ & $330 = 66 \cdot 5$  & $66$  \\
$2$ & $120 = 30 \cdot 4$  & $30$  \\
$3$ & $36 = 12 \cdot 3$   & $12$  \\
$4$ & $8 = 4 \cdot 2$     & $4$   \\
$5$ & $1 = 1 \cdot 1$     & $1$   \\
\bottomrule
\end{tabular}
\end{center}

\begin{remark}
Thus the six conditions hold simultaneously at $n = 12$. There is \emph{no}
divisibility obstruction at $n = 12$, in direct contradiction to the claim that
the bound $n_0(5,6) > 12$ is supplied by such an obstruction.
\end{remark}

\section{The existence certificate}

\begin{definition}
A $t$-$(n,k,1)$ design (Steiner system $S(t,k,n)$) is a family $\mathcal{B}$ of
$k$-subsets (blocks) of an $n$-set such that every $t$-subset lies in exactly
one block. For $t=5$, $k=6$, $n=12$ the number of blocks must be
$\lambda_0 = \binom{12}{5}/\binom{6}{5} = 132$.
\end{definition}

\begin{theorem}\label{thm:existence}
There exists a $5$-$(12,6,1)$ design: $132$ six-element blocks on the point set
$\{0,1,\dots,11\}$ such that every one of the $\binom{12}{5} = 792$ five-element
subsets is contained in exactly one block.
\end{theorem}

\begin{proof}[Two independent constructions, each exhaustively verified]
We give two constructions and check both against all $792$ five-subsets.

\emph{(a) Exact cover / backtracking.} Let the columns be the $792$
five-subsets and the rows the $924 = \binom{12}{6}$ six-subsets; selecting a row
covers its $\binom{6}{5} = 6$ five-subsets. An Algorithm~X backtracking search
with the minimum-remaining-values heuristic over the $792$ columns selects $132$
rows. Direct enumeration confirms that all $792$ five-subsets are covered
\emph{exactly once} (each count equals $1$), that the $132$ blocks are distinct
and all have size $6$. This is the design whose $132$ $12$-bit masks appear
verbatim in \texttt{lean4/Main.lean}.

\emph{(b) Extended ternary Golay code.} The extended ternary Golay code
$[12,6,6]$ over $\mathbb{F}_3$ has weight distribution
\[
  A_0 = 1,\quad A_6 = 264,\quad A_9 = 440,\quad A_{12} = 24
\]
and $3^6 = 729$ codewords. Its $264$ weight-$6$ words split into $132$ opposite
pairs $\{\pm c\}$, giving $132$ distinct $6$-element supports. Exhaustive
enumeration of all $792$ five-subsets confirms that these $132$ supports also
cover every five-subset exactly once.

\emph{The two designs are isomorphic.} An explicit point permutation
$p = (0,1,2,3,4,11,7,9,8,6,5,10)$ maps the exact-cover design onto the Golay
design; i.e.\ the two labelled designs coincide up to relabelling. The common
design is $\Wtwelve$.
\end{proof}

\begin{remark}[Invariants and uniqueness]
Both constructions satisfy the same combinatorial invariants: the
$132 \cdot 131 / 2 = 8646$ pairs of blocks intersect in
$0,2,3,4$ points with multiplicities $66, 2970, 2640, 2970$ respectively
(intersections of size $1$ or $5$ never occur), and each block meets the others
in $i$ points with multiplicities $1,0,45,40,45,0,1$ for $i = 0,\dots,6$. The
design is classically unique up to isomorphism, with automorphism group
$M_{12}$ of order $95040$; \texttt{reproduce.py} confirms this computationally by
enumerating all $95040$ automorphisms of the constructed design.
\end{remark}

\section{The logical dichotomy}

The conjecture is a conjunction, and no reading makes both clauses true.

\begin{theorem}\label{thm:dichotomy}
Let $n_0(5,6)$ be the threshold appearing in clause (C1). Then:
\begin{enumerate}
\item[(i)] if (C1) holds, then $n_0(5,6) \le 12$, so (C2) is false;
\item[(ii)] if (C2) holds, then (C1) is false.
\end{enumerate}
Hence (C1) $\wedge$ (C2) is false under every reading.
\end{theorem}

\begin{proof}
By Theorem~\ref{thm:existence} there is an $S(5,6,12)$, and by
Section~2 the six divisibility conditions hold at $n=12$.

(i) Assume (C1). Instantiating it at $(t,k,n) = (5,6,12)$ gives
\[
  \bigl(\text{$S(5,6,12)$ exists}\bigr)
  \iff
  \bigl(\text{divisibility holds at } (5,6,12) \;\wedge\; 12 \ge n_0(5,6)\bigr).
\]
The left side is true (Theorem~\ref{thm:existence}) and the divisibility
conjunct on the right is true (Section~2), so the equivalence forces
$12 \ge n_0(5,6)$, i.e.\ $n_0(5,6) \le 12$. The literal statement of (C2),
$n_0(5,6) > 12$, is therefore false, and so is the claim that this bound is
``given by a concrete divisibility obstruction'' (there is no obstruction at
$n = 12$: all six conditions hold).

(ii) Assume (C2), i.e.\ $n_0(5,6) > 12$. At $(t,k,n) = (5,6,12)$ the
divisibility conditions hold (Section~2) but $12 \ge n_0(5,6)$ fails, so the
right side of (C1) is false, while the left side is true by
Theorem~\ref{thm:existence}. Thus the asserted ``if and only if'' in (C1) is
false.

Either way the conjunction (C1) $\wedge$ (C2) fails. This holds whether
$n_0(t,k)$ is read as the exact least admissible $n$, as any threshold function
making (C1) true, or as the explicit bound of type $(kt)^{ct}$ mentioned in the
conjecture: any such quantity is forced by (C1) to satisfy
$n_0(5,6) \le 12$.
\end{proof}

\begin{corollary}
Conjecture \texttt{00000008417} is FALSE.
\end{corollary}

\begin{proof}
The conjecture asserts (C1) $\wedge$ (C2), which Theorem~\ref{thm:dichotomy}
shows to be unsatisfiable. Clause (C2) is independently and concretely false:
the divisibility conditions hold at $n=12$ and an $S(5,6,12)$ exists, so the
threshold is at most $12$.
\end{proof}

\section{Threshold context}

In every small case the first admissible proper value of $n$ is attained, and
the extremal examples are exactly the Witt designs.

\begin{center}
\begin{tabular}{c c l l}
\toprule
$(t,k)$ & first admissible $n$ & design & name \\
\midrule
$(2,3)$ & $7$  & $S(2,3,7) = \mathrm{PG}(2,2)$ & Fano plane \\
$(3,4)$ & $8$  & $S(3,4,8) = \mathrm{AG}(3,2)$ & affine geometry \\
$(4,5)$ & $11$ & $S(4,5,11)$ & Witt design $W_{11}$ \\
$(5,6)$ & $12$ & $S(5,6,12)$ & Witt design $\Wtwelve$ \\
$(5,8)$ & $24$ & $S(5,8,24)$ & Witt design $W_{24}$ \\
\bottomrule
\end{tabular}
\end{center}

The $(5,6)$ row is decisive: divisibility holds at $n = 12$ and the design
exists there, so the threshold satisfies $n_0(5,6) \le 12 < n_0(5,6)$ under the
filed bound.

\section{Scope note}

We are explicit about what is and is not in dispute.

\begin{itemize}
\item The general large-$n$ criterion ``divisibility implies existence for
  $n \ge n_0(t,k)$'' is essentially Keevash's theorem. That theorem is
  non-effective and there is no explicit $(kt)^{ct}$ threshold in the
  literature. \emph{This submission does not dispute the large-$n$ existence
  direction.} It targets the definite, false conjunct $n_0(5,6) > 12$ (and the
  accompanying ``concrete divisibility obstruction'').
\item At $n = 12$ the divisibility conditions hold and the design exists; hence
  no lower bound above $12$ can be correct for $(5,6)$, and no divisibility
  obstruction at $n = 12$ exists. This is a finite, verified, unconditional
  refutation.
\item Uniqueness of $\Wtwelve$ up to isomorphism and its automorphism group
  $M_{12}$ (order $95040$) are classical; they are not needed for the
  refutation, which needs only one design. They are reported for context and
  checked computationally.
\end{itemize}

\section{Formalisation}

The existence certificate is formalised in \texttt{lean4/} (core Lean only,
\texttt{import Std}, no Mathlib, no \texttt{sorry}). Points are encoded as
$12$-bit \texttt{Nat} masks. The central theorem is

\begin{quote}\ttfamily
theorem every\_five\_in\_exactly\_one\_block :\\
\hspace*{2em}fives.all (fun m => covered m == 1) = true
\end{quote}

proved by \texttt{decide}; it scans all $792$ five-subsets against the $132$
blocks of \texttt{blocks} and checks that each is contained in exactly one. The
six divisibility identities (\texttt{div\_i0}, \dots, \texttt{div\_i5}),
\texttt{witt\_is\_design} (\texttt{IsDesign5612 blocks}) and
\texttt{exists\_S5612} ($\exists B$, $B$ is a $5$-$(12,6,1)$ design) are also
proved. The audit \texttt{lake env lean Check.lean} reports only the axiom
\texttt{propext} --- no \texttt{sorryAx}, and no \texttt{Lean.ofReduceBool}
because the fully kernel-checked \texttt{decide} variant is used. See
\texttt{lean4/README.md} for the build time ($\approx 63$~s) and the full
statement table.

\section{Reproducibility}

\texttt{reproduce.py} uses the Python~3 standard library only and performs:
(a) the six divisibility quotients, checked to be integers;
(b) the exact-cover construction with an exhaustive check that every one of the
$792$ five-subsets lies in exactly one of the $132$ blocks, plus a SHA-256 of
the canonical block list;
(c) the independent extended ternary Golay construction with the same
exhaustive coverage check;
(d) an explicit isomorphism between the two labelled designs.

\begin{quote}\ttfamily
\$ python3 reproduce.py\\
...\\
PASS: all checks verified; conjecture 00000008417 is FALSE.
\end{quote}

It prints the block count, the canonical SHA-256
(\texttt{6b2ab5e7\dots eb567}), the invariants, an explicit relabelling
isomorphism, and $|\mathrm{Aut}| = 95040$, and exits $0$ exactly when every
check holds.

\begin{quote}\ttfamily
\$ tectonic --outdir build main.tex\\
\$ ls -l build/main.pdf
\end{quote}

\section*{Status against the submission rules}

Rule 3 requires a LaTeX source, a PDF document, and a Lean~4 project.

\begin{itemize}
\item \textbf{LaTeX source} --- \texttt{main.tex}, a standalone \texttt{article}
  using \texttt{geometry}, \texttt{amsmath}, \texttt{amssymb}, \texttt{amsthm},
  \texttt{array}, \texttt{parskip}, \texttt{booktabs}; compiles with
  \texttt{tectonic main.tex}.
\item \textbf{PDF} --- at \texttt{build/main.pdf}, produced by
  \texttt{tectonic --outdir build main.tex}.
\item \textbf{Lean~4 project} --- \texttt{lean4/} with \texttt{lean-toolchain}
  pinned to \texttt{leanprover/lean4:v4.33.1}, \texttt{lakefile.toml}
  (library \texttt{Main}, \texttt{name = "tlmc8417"}, no dependencies),
  \texttt{Main.lean}, \texttt{Check.lean} (\texttt{\#print axioms}), and
  \texttt{lean4/README.md}. \texttt{lake build} exits $0$ and the audit reports
  no \texttt{sorryAx}.
\end{itemize}

\end{document}
